% L2b takeaway deck.
%
% THE DECK MIRRORS THE NOTEBOOK. The instructor shows
% CHEME-5660-L2b-Lecture-Yield-Duration-Convexity-Curve-Fall-2026.ipynb while
% students follow along and take notes on these slides. The deck therefore
% covers the same material in the same order, with frame titles that state the
% main point of each notebook section. Adding a section to the notebook means
% adding a slide here.
% See the content conventions in templates/vnslides.sty.
%
% Notebook outline, and where each part lands:
%
%   (deck only)                                  -> title
%   Disclaimer and Risks              [last]     -> slide 2  (always near the front)
%   Motivation: $500M scenario                   -> Scenario
%   Learning objectives               [cell 0]   -> Objectives for Today
%   Examples                                     -> Lectures and Examples
%   Review: default and interest-rate risk       -> transition frame + review slide
%     SVB discussion                             -> SVB: Rate Risk Met Deposit Flight
%   Motivation: Malkiel's theorems stated        -> Malkiel's Bond-Pricing Theorems
%   Zero-Coupon Treasury Bill                    -> A Treasury Bill Has One Cash Flow
%     Price Sensitivity of a Treasury Bill       -> A Bill's Price Depends on Yield and Time
%     Interpreting the Bill Sensitivities        -> What Does a Change in Maturity Mean?
%                                                   + Yield and Maturity Sensitivities
%                                                   + Higher Yield Lowers Price ...
%     Optional zero-coupon extension             -> Optional: Bill Maturity Sensitivity
%   Coupon-Bearing Treasury Notes and Bonds      -> cash-flow + sensitivity slides
%     Interpreting the Coupon-Bond Sensitivities -> How Bond Features Change Its Price
%   Duration and Convexity                       -> common price formula + derivatives
%     Price-Yield Geometry                       -> curve shape + figure
%     Duration and Convexity as Risk Measures    -> normalized measures + Macaulay comparison
%     Estimating a Price Change                  -> approximation + worked estimate + portfolio use
%   Revisiting Malkiel's Theorems                -> par duration + four theorem slides
%   From One Yield to the Yield Curve            -> one yield versus maturity-specific rates
%     recession-predictor blockquote             -> inverted-curve warning
%   Optional Advanced Material                   -> Optional Advanced Material
%   Summary                                      -> What to Remember
%   (deck only)                                  -> closing transition frame
%
% Like L2a and unlike L1b, this notebook DOES open with a concept-review section,
% so \transitionframe is used as intended rather than as a 2024 holdover.
%
% Figure: L2b has no figs/ directory. The price-yield convexity figure is
% inlined here and scaled up
% for projection. It carries the lecture's central geometry, so the deck should
% not be text-only just because the notebook has no figure. The palette names it
% uses (vnink, vnlink, vncarnelian, vnrule, vnmuted) reach this file through
% templates/vnflow.sty, which vnslides.sty inputs.
%
% NOTATION. The deck and notebook agree and follow the
% course-wide conventions set in L1b and L2a:
%   price            V_B            (the notes' \mathcal{B} was retired)
%   coupon rate      c              (L2b's \bar c was retired; L2a uses c)
%   discount factor  \mathcal D_{j,0}^{-1}(y;n)   -- note the ;n argument
%   bill factor      \mathcal D_{T,0}^{-1}(y;n) = (1+y/n)^{-nT}
%   continuous rate  g_y := n\log(1+y/n)
% There is deliberately NO one-letter growth factor. L2b used to write q := 1+y/n,
% which collided with the risk-neutral up probability q in L3b and in this
% lecture's own advanced/yield_curve notes. Per commit ff4dbb3, q is spoken for.
% Discount factors are named \mathcal D^{-1} in valuation statements, but written
% out as (1+y/n)^{-j} wherever a derivative is being taken, so the power rule
% stays visible.
%
% The notebook's V_{B,y} and V_{B,T} are kept for the two sensitivity MAGNITUDES.
% A reviewer flagged the comma subscript as looking like a partial derivative; it
% was kept anyway, because diverging from the companion notebook is the worse
% fault under the governing rule. The frame states explicitly that the
% derivatives are their negatives.
\documentclass[aspectratio=169,11pt]{beamer}
\usepackage{vnslides}

\title{{\vnheavy L2b}: Yield, Duration, Convexity, and the Yield Curve}
\subtitle{CHEME 5660 Cornell University}
\author{Copyright Jeffrey Varner 2026. All Rights Reserved}

\begin{document}

% ---- title ---------------------------------------------------
\vntitlepage

% ---- disclaimer ----------------------------------------------
\begin{frame}{Disclaimer and Risks}
  \vspace{1.0em}
  \textbf{\ink{This material is for training and information only.}} It does not
  provide investment advice or recommend any security, derivative, trade, or
  strategy.

  \vspace{1.2em}
  \textbf{\ink{Trading involves risk.}} Past results do not guarantee future
  results. Any risk you take should fit your resources, experience, and
  tolerance for loss.

  \vspace{1.2em}
  \textbf{\ink{Your decisions are your responsibility.}} Base them on your own
  objectives, financial position, risk tolerance, and liquidity needs.
\end{frame}

% ---- scenario ------------------------------------------------
\begin{frame}{Scenario}
  \vspace{0.4em}
  Suppose you manage a \hilite{\$500 million Treasury portfolio}. The Federal
  Reserve meets today, and the market expects a possible 50-basis-point move in
  the policy rate.

  \vspace{1.0em}
  If yields move, the value of the portfolio will move with them, and your
  clients will want to know \hilite{by how much}.

  \vspace{1.0em}
  The estimate depends on the timing of the portfolio's cash flows and on the
  change in each security's yield. We will use the \hilite{slope} and
  \hilite{curvature} of the price-yield curve to compute the price response.

  \slidenote{Policy rate and yield:}{A half-point move in the \emph{policy rate} does not imply a half-point move in \emph{your security's yield}.}
\end{frame}

% ---- objectives ----------------------------------------------
\begin{frame}{Learning Objectives}
  \vspace{0.4em}
  We start with the L2a pricing model and study how fixed-income prices respond
  when yields change.

  \vspace{0.9em}
  {\large\textbf{\ink{Today's concepts}}}
  \vspace{0.4em}
  \begin{itemize}\setlength{\itemsep}{0.5em}
    \item \hilite{Price and yield}: Explain why a higher yield lowers the value of
      fixed positive cash flows and why the price-yield curve is convex.
    \item \hilite{Duration and convexity}: Use both measures to estimate a price
      change for a small change in yield, and state the units of each measure.
    \item \hilite{Yield measures}: Distinguish yield to maturity, par yield, spot
      rates, and maturity-specific discount factors.
  \end{itemize}
  \vspace{-0.4em}

  \slidenote{Companion notebook:}{\href{https://github.com/varnerlab/CHEME-5660-CourseRepository-Fall-2026}{L2b lecture notebook}}
\end{frame}

% ---- examples ------------------------------------------------
\begin{frame}{Lecture and Example Notebooks}
  \vspace{0.5em}
  {\normalsize\textbf{\ink{Lecture:}} \dimmed{CHEME-5660-L2b-Lecture-Yield-Duration-Convexity-Curve-Fall-2026.ipynb}}

  \vspace{1.0em}
  {\normalsize\textbf{\ink{Example:}} \dimmed{CHEME-5660-L2b-Example-Treasury-Duration-Convexity-Fall-2026.ipynb}}

  \vspace{1.0em}
  The example prices a coupon Treasury security and computes the price response
  to changes in yield and coupon rate.

  \vspace{0.8em}
  The \dimmed{advanced/} directory contains optional notebooks on Malkiel's
  theorems and the yield curve.

  \slidenote{Link:}{\href{https://github.com/varnerlab/CHEME-5660-CourseRepository-Fall-2026.git}{https://github.com/varnerlab/CHEME-5660-CourseRepository-Fall-2026.git}}
\end{frame}

% ---- concept review ------------------------------------------
\transitionframe{Treasury prices change when \hilite{market yields} change.}

\begin{frame}{Are Treasury Securities Risk-Free?}
  \vspace{0.3em}
  \hilite{Not exactly.} U.S. Treasury securities have low default risk, but they
  are still exposed to interest-rate risk.

  \vspace{0.45em}
  \begin{itemize}\setlength{\itemsep}{0.45em}
    \item \hilite{Default risk}: The issuer fails to make a promised payment.
      This risk is low for U.S. Treasury securities, although sovereign defaults
      do occur.
    \item \hilite{Interest-rate risk}: A change in market yields changes the
      value of fixed promised cash flows. The market value can fall even if the
      security is not sold.
  \end{itemize}

  \slidenote{Focus:}{We focus on interest-rate risk because we can compute how a change in yield changes the price.}
\end{frame}

\begin{frame}{Silicon Valley Bank and Interest-Rate Risk}
  \vspace{0.1em}
  Silicon Valley Bank held many \hilite{long-maturity, fixed-rate} government
  securities. Their market values fell when interest rates rose in 2022 and
  2023.

  \vspace{0.3em}
  \begin{itemize}\setlength{\itemsep}{0.3em}
    \item Many of the losses remained unrealized while the bank held the
      securities.
    \item In March 2023, uninsured depositors withdrew funds. To meet those
      withdrawals, SVB sold securities, realized losses, and announced a capital
      raise.
    \item The FDIC took control on March 10, 2023, when the bank had roughly
      \$209 billion in assets. At the time, it was the second-largest failure of
      an insured bank on record; it is now the third-largest.
  \end{itemize}

  \vspace{0.15em}
  The decline in asset values and the deposit withdrawals reinforced each
  other. Interest-rate risk was important, but it was not the only cause of the
  failure.

  \slidenote{Overview:}{\href{https://www.youtube.com/watch?v=7tNt7WaP4iA}{the collapse of Silicon Valley Bank, from PBS}}
\end{frame}

% ---- Malkiel theorems stated up front ------------------------
\begin{frame}{Malkiel's Bond-Pricing Theorems}
  \vspace{0.2em}
  Malkiel's five theorems describe the prices of \hilite{option-free} bonds with
  fixed positive cash flows. Each comparison holds the other contract terms
  fixed. The maturity comparisons use bonds priced \hilite{at par}, so $c=y$.

  \vspace{0.4em}
  \begin{enumerate}\setlength{\itemsep}{0.35em}
    \item The price falls when the yield rises.
    \item A longer-maturity bond is more sensitive to a change in yield.
    \item Each additional period adds less sensitivity than the previous one.
    \item A decrease in yield raises the price more than an equal increase in
      yield lowers it.
    \item A higher coupon rate reduces the percentage price response to a given
      change in yield.
  \end{enumerate}

  \slidenote{Reference:}{Malkiel (1962), \emph{Expectations, Bond Prices, and the Term Structure of Interest Rates}, QJE 76(2), 197\textendash218.}
\end{frame}

% ---- zero-coupon bill ----------------------------------------
\begin{frame}{Treasury Bill Pricing}
  \vspace{0.3em}
  A Treasury bill has one cash flow: the face value $V_P$ paid at maturity.

  \vspace{0.5em}
  Let $T$ be the remaining time to maturity in years, $n$ the number of
  compounding periods per year, and $y$ the nominal annual yield. From L2a, the
  price is:
  \[
    V_B=\underbrace{\left(1+y/n\right)^{-nT}}_{\mathcal D_{T,0}^{-1}(y;n)}V_P .
  \]

  \vspace{0.4em}
  A larger yield or a longer maturity discounts the face value more heavily. We
  assume $V_P>0$, $T>0$, and $y>0$ unless stated otherwise.
\end{frame}

\begin{frame}{Price Sensitivity of a Treasury Bill}
  \vspace{0.2em}
  A bill has no coupon, so $c=0$ and $dc=0$. Holding $V_P$ and $n$ fixed, its
  price changes with the yield $y$ and the remaining maturity $T$:
  \[
    dV_B
    =-\underbrace{V_P\left(1+y/n\right)^{-nT}\left(\frac{T}{1+y/n}\right)}_{V_{B,y}}dy
    \;-\;\underbrace{V_P\left(1+y/n\right)^{-nT}\bigl(n\log\left(1+y/n\right)\bigr)}_{V_{B,T}}dT .
  \]

  \vspace{0.4em}
  The quantities $V_{B,y}$ and $V_{B,T}$ are positive sensitivity magnitudes.
  The minus signs show that the price falls when either the yield or the
  remaining maturity increases, provided $y>0$.

  \slidenote{No $dc$ term:}{The contract fixes $dc=0$. The derivative $\partial V_B/\partial c$ is not zero, even at $c=0$.}
\end{frame}

\begin{frame}{Interpreting a Change in Maturity}
  \vspace{0.3em}
  A change in $T$ can describe either a comparison between contracts or the
  passage of time for one contract:

  \vspace{0.6em}
  \begin{itemize}\setlength{\itemsep}{0.55em}
    \item \hilite{Compare two contracts}: A 6-month bill and a 12-month bill have
      different contractual maturities. For a coupon bond, a longer maturity can
      also add coupon payments.
    \item \hilite{Hold one contract over time}: As calendar time advances, the
      remaining maturity falls, so $dT=-dt$. For a coupon bond, this change is
      smooth only between coupon dates.
  \end{itemize}

\end{frame}

\begin{frame}{Treasury Bill Sensitivities}
  \vspace{-0.3em}
  The sensitivity magnitudes for yield and maturity are:
  \[
    V_{B,y}=\underbrace{V_P(1+y/n)^{-nT}}_{=\,V_B}\underbrace{\left(\frac{T}{1+y/n}\right)}_{>\,0},
    \qquad
    V_{B,T}=\underbrace{V_P(1+y/n)^{-nT}}_{=\,V_B}\underbrace{\bigl(n\log(1+y/n)\bigr)}_{=\,g_y},
  \]
  where $g_y:=n\log(1+y/n)$ is the continuously compounded rate equivalent to
  $y$.

  \vspace{0.3em}
  \begin{itemize}\setlength{\itemsep}{0.25em}
    \item $V_{B,y}>0$ for $T>0$ and $1+y/n>0$.
    \item $V_{B,T}>0$ requires $y>0$.
    \item At $y=0$, maturity does not affect price. For $-n<y<0$, a longer bill
      has a higher price.
  \end{itemize}

\end{frame}

\begin{frame}{How Yield and Time Affect a Treasury Bill}
  \vspace{0.4em}
  {\large\textbf{\ink{Effect of a higher yield}}}
  \vspace{0.2em}

  A higher yield discounts the face value more heavily, so
  $y\uparrow \Rightarrow V_B\downarrow$. If a 52-week bill's yield rises from
  2\% to 3\%, its market price falls.

  \vspace{0.9em}
  {\large\textbf{\ink{Effect of time passing}}}
  \vspace{0.2em}

  At a fixed positive yield, the price rises as the bill approaches maturity
  because fewer discounting periods remain: $T\downarrow \Rightarrow
  V_B\uparrow$. At $y=0$, the discount factor is one, so the price does not
  depend on maturity.

\end{frame}

\begin{frame}{Optional Treasury Bill Example}
  \vspace{0.6em}
  The required lecture continues with coupon-bearing notes and bonds. The
  optional bill notebook examines maturity sensitivity in more detail.

  \vspace{1.0em}
  {\normalsize\textbf{Question:} How does the price sensitivity of a zero-coupon
  bill change with its remaining maturity?}

  \vspace{1.0em}
  {\normalsize\textbf{Notebook:} The Treasury bill maturity example in
  \dimmed{advanced/malkiel/}.}
\end{frame}

% ---- coupon notes and bonds ----------------------------------
\begin{frame}{Coupon-Bearing Treasury Notes and Bonds}
  \vspace{0.1em}
  A coupon-bearing Treasury security pays coupons before maturity and returns
  the face value at maturity. Let $c$ be the annual coupon rate,
  $C:=(c/n)V_P$ the coupon paid each period, and $N:=nT$ the number of remaining
  payment dates. The L2a pricing model is:
  \[
    V_B=\mathcal D_{N,0}^{-1}(y;n)\,V_P+C\sum_{j=1}^{N}\mathcal D_{j,0}^{-1}(y;n),
    \qquad \mathcal D_{j,0}^{-1}(y;n)=\left(1+y/n\right)^{-j}.
  \]

  \vspace{0.5em}
  The price has two parts:
  \begin{itemize}\setlength{\itemsep}{0.4em}
    \item the discounted principal payment $V_P$ at date $N$;
    \item the discounted coupon payments $C$ at dates $j=1,\ldots,N$.
  \end{itemize}
\end{frame}

\begin{frame}{Yield and Coupon Sensitivity}
  \vspace{-0.2em}
  Hold the payment schedule fixed and differentiate:
  \[
    \begin{aligned}
      \frac{\partial V_B}{\partial y}
        &=-\frac{1}{n(1+y/n)}\left[NV_P\,\mathcal D_{N,0}^{-1}(y;n)
          +C\sum_{j=1}^{N}j\,\mathcal D_{j,0}^{-1}(y;n)\right]<0,\\[0.2em]
      \frac{\partial V_B}{\partial c}
        &=\frac{V_P}{n}\sum_{j=1}^{N}\mathcal D_{j,0}^{-1}(y;n)>0 .
    \end{aligned}
  \]

  \vspace{0.1em}
  A higher yield discounts every payment more heavily, so the price falls. A
  higher coupon rate increases every coupon payment, so the price rises.

  \slidenote{Schedule length:}{$N$ is a discrete index. We do not differentiate a finite sum with upper limit $nT$ as if $T$ were continuous.}
\end{frame}

\begin{frame}{How Yield, Coupon, and Maturity Affect Price}
  \vspace{0.3em}
  \begin{itemize}\setlength{\itemsep}{0.55em}
    \item \hilite{Yield}: $y\uparrow \Rightarrow V_B\downarrow$ because every
      coupon and the principal payment are discounted more heavily.
    \item \hilite{Coupon rate}: $c\uparrow \Rightarrow V_B\uparrow$, but not in
      direct proportion. Doubling $c$ doubles the coupon stream, not the
      principal payment.
    \item \hilite{Remaining maturity}: The direction of the price change depends
      on $c-y$. As maturity shortens, a discount bond ($c<y$) rises toward par, a
      premium bond ($c>y$) falls toward par, and a par bond ($c=y$) stays at par
      on coupon dates.
  \end{itemize}

  \slidenote{Maturity comparison:}{Malkiel's result concerns \emph{sensitivity to a yield change}, not the price level.}
\end{frame}

% ---- duration and convexity ----------------------------------
\begin{frame}{A General Cash-Flow Pricing Model}
  \vspace{0.2em}
  For an \hilite{option-free} security, the promised cash flows do not change
  when yields change. Conventional noncallable Treasury securities satisfy this
  assumption.

  \vspace{0.5em}
  Let $X_j\ge0$ be the cash flow at period $j=1,\ldots,N$, with at least one
  positive cash flow. The price is:
  \[
    V_B(y)=\sum_{j=1}^{N}X_j\,\mathcal D_{j,0}^{-1}(y;n).
  \]

  \vspace{0.3em}
  \begin{itemize}\setlength{\itemsep}{0.3em}
    \item \hilite{Note or bond}: $X_j=C$ for $j<N$ and $X_N=C+V_P$.
    \item \hilite{Bill}: $X_j=0$ for $j<N$ and $X_N=V_P$.
  \end{itemize}

  \slidenote{Outside this model:}{callable, inflation-indexed, and mortgage-backed securities can have cash flows that change with rates or states.}
\end{frame}

\begin{frame}{The Price-Yield Relationship}
  \vspace{0.3em}
  Holding the cash flows $X_j$ and the number of payments $N$ fixed gives:
  \[
    \begin{aligned}
      \frac{dV_B}{dy}&=-\frac{1}{n}\sum_{j=1}^{N}j\,X_j\left(1+y/n\right)^{-j-1}<0,\\[0.2em]
      \frac{d^2V_B}{dy^2}&=\frac{1}{n^2}\sum_{j=1}^{N}j(j+1)X_j\left(1+y/n\right)^{-j-2}>0 .
    \end{aligned}
  \]

  \vspace{0.4em}
  \begin{itemize}\setlength{\itemsep}{0.35em}
    \item $dV_B/dy<0$: The price decreases as the yield increases.
    \item $d^2V_B/dy^2>0$: The price-yield curve is strictly convex.
  \end{itemize}

  \slidenote{Sign conditions:}{$1+y/n>0$, every $X_j\ge0$, and at least one $X_j>0$.}
\end{frame}

\begin{frame}{Convexity and Asymmetric Price Changes}
  \vspace{-0.4em}
  \begin{center}
    \begin{tikzpicture}[x=1.28cm,y=0.84cm,>=Latex,font=\sffamily\small]
      % axes
      \draw[vnink,line width=0.8pt,-{Latex[length=3mm]}] (0,0)--(7.2,0) node[right]{yield $y$};
      \draw[vnink,line width=0.8pt,-{Latex[length=3mm]}] (0,0)--(0,4.7) node[above]{price $V_B(y)$};
      % tangent at y0: the duration (straight-line) estimate
      \draw[vnmuted,line width=0.9pt,domain=1.2:6.05]
        plot (\x,{2.388-0.456*(\x-3.45)});
      \node[anchor=west,text=vnmuted,align=left] at (6.1,1.0)
        {\footnotesize tangent at $y_0$:\\[-2pt]\footnotesize the duration line};
      % price curve (domain chosen so the curve stays inside the axes)
      \draw[vnlink,line width=1.4pt,domain=0.96:6.7,samples=80,smooth]
        plot (\x,{5.25/(0.75+0.42*\x)});
      % the three priced points
      \coordinate (mid)  at (3.45,2.388);
      \coordinate (low)  at (1.55,3.748);
      \coordinate (high) at (5.35,1.752);
      % droplines to the yield axis and guides to the price axis
      \draw[vnrule,dashed] (1.55,0)--(low) (3.45,0)--(mid) (5.35,0)--(high);
      \draw[vnrule,dashed] (0,3.748)--(low) (0,2.388)--(mid) (0,1.752)--(high);
      \fill[vnink] (mid) circle (1.8pt);
      \fill[vnlink] (low) circle (1.8pt);
      \fill[vncarnelian] (high) circle (1.8pt);
      \node[below] at (1.55,0) {$y_-$};
      \node[below] at (3.45,0) {$y_0$};
      \node[below] at (5.35,0) {$y_+$};
      % equal shocks marked below the axis
      \draw[vnmuted,<->,line width=0.8pt] (1.55,-0.62)--(3.45,-0.62)
        node[midway,below,text=vnmuted]{\footnotesize $\Delta y$};
      \draw[vnmuted,<->,line width=0.8pt] (3.45,-0.62)--(5.35,-0.62)
        node[midway,below,text=vnmuted]{\footnotesize $\Delta y$};
      % gain and loss stacked on the price axis, sharing the V(y0) endpoint
      \draw[vnlink,<->,line width=1.0pt] (-0.22,2.388)--(-0.22,3.748)
        node[midway,left,text=vnlink]{gain};
      \draw[vncarnelian,<->,line width=1.0pt] (-0.22,1.752)--(-0.22,2.388)
        node[midway,left,text=vncarnelian]{loss};
    \end{tikzpicture}
  \end{center}

  \vspace{-0.5em}
  {\footnotesize For equal changes in yield, the \hilite{gain} from a decrease is
  larger than the \hilite{loss} from an increase. The tangent line gives the
  duration estimate, while convexity accounts for the curvature.}
\end{frame}

\begin{frame}{Modified Duration and Convexity}
  \vspace{0.2em}
  The price derivatives measure dollar sensitivity, so they depend on the size
  of the position. Dividing by the current price $V_B(y)>0$ gives percentage
  sensitivity measures that can be compared across securities:
  \[
    D_{\mathrm{mod}}(y):=-\frac{1}{V_B(y)}\frac{dV_B}{dy},
    \qquad
    K(y):=\frac{1}{V_B(y)}\frac{d^2V_B}{dy^2} .
  \]
  \begin{itemize}\setlength{\itemsep}{0.35em}
    \item \hilite{Modified duration $D_{\mathrm{mod}}$} measures the local
      percentage slope of price with respect to yield.
    \item \hilite{Convexity $K$} measures the local percentage curvature.
    \item If $y$ is a decimal annual yield, $D_{\mathrm{mod}}$ has units of years
      and $K$ has units of years squared.
  \end{itemize}
\end{frame}

\begin{frame}{Macaulay and Modified Duration}
  \vspace{0.5em}
  \hilite{Macaulay duration} is the present-value-weighted average time to
  payment. \hilite{Modified duration} converts that measure of cash-flow timing
  into local price sensitivity:
  \[
    D_{\mathrm{mod}}=\frac{D_{\mathrm{mac}}}{1+y/n}.
  \]

  \vspace{0.7em}
  A bill makes its only payment at maturity, so:
  \[
    D_{\mathrm{mac}}=T,
    \qquad
    D_{\mathrm{mod}}=\frac{T}{1+y/n}=\frac{V_{B,y}}{V_B}.
  \]
\end{frame}

\begin{frame}{Estimating a Price Change}
  \vspace{0.2em}
  For a small yield change $\Delta y$, the second-order Taylor approximation is:
  \[
    \frac{\Delta V_B}{V_B}\approx
    \underbrace{-D_{\mathrm{mod}}\,\Delta y}_{\text{\footnotesize slope}}
    +\underbrace{\tfrac{1}{2}K(\Delta y)^2}_{\text{\footnotesize curvature}} .
  \]
  Because $K>0$, the convexity term is positive for either sign of $\Delta y$.
  It increases the estimated gain when the yield falls and reduces the estimated
  loss when the yield rises.

  \vspace{0.4em}
  \hilite{Use decimal yield changes:} one basis point is $10^{-4}$, so 50 basis
  points is $\Delta y=0.0050$.

  \slidenote{Policy rate and yield:}{A 50-basis-point move in the \emph{policy rate} does not imply $\Delta y=0.0050$ for a given security.}
\end{frame}

\begin{frame}{Example: A 50-Basis-Point Yield Increase}
  \vspace{0.4em}
  Suppose a note has $D_{\mathrm{mod}}=7.2$ years and $K=65$ years$^2$. If its
  yield rises by 50 basis points, then $\Delta y=0.0050$:

  \vspace{0.5em}
  \begin{itemize}\setlength{\itemsep}{0.4em}
    \item The duration term is $-7.2\times0.0050=-0.0360$.
    \item The convexity term is
      $\tfrac{1}{2}\times65\times(0.0050)^2=+0.0008$.
    \item The estimated price change is $\hilite{-3.52\%}$, compared with
      $-3.60\%$ using duration alone.
  \end{itemize}

  \vspace{0.6em}
  Duration and convexity are \hilite{local} measures. For a large yield change,
  reprice the full cash-flow schedule at the new yield.
\end{frame}

\begin{frame}{Portfolio Duration}
  \vspace{0.3em}
  \begin{itemize}\setlength{\itemsep}{0.55em}
    \item Duration places securities with different coupons, maturities, and
      prices on the same sensitivity scale.
    \item Portfolio duration is the market-value-weighted average of the
      durations of the individual positions.
    \item Offsetting positions can reduce net duration. Funding long-duration
      assets with short-duration liabilities creates a duration mismatch.
  \end{itemize}

  \vspace{0.5em}
  These calculations assume that each position experiences the stated yield
  change. If different maturities move by different amounts, we need
  maturity-specific yield changes.
\end{frame}

% ---- Malkiel -------------------------------------------------
\begin{frame}{Duration of a Par Bond}
  \vspace{0.2em}
  To compare sensitivity across maturities, consider option-free bonds with
  fixed positive cash flows and hold the other terms fixed. Each bond starts at
  \hilite{par}, so $c=y$.

  \vspace{0.4em}
  Sum the coupon-bond price, differentiate at fixed $c$, and then set $c=y$:
  \[
    \frac{V_B}{V_P}=\frac{c}{y}+\left(1-\frac{c}{y}\right)\left(1+y/n\right)^{-N}
    \qquad\Longrightarrow\qquad
    D_{\mathrm{mod}}(N)=\frac{1-\left(1+y/n\right)^{-N}}{y} .
  \]

  \vspace{0.4em}
  The order of these steps matters. If we set $c=y$ before differentiating, the
  price remains at $V_P$ along the par curve and the derivative is zero. We
  assume $y>0$.
\end{frame}

\begin{frame}{Malkiel's Theorems 1 and 2}
  \vspace{0.4em}
  {\large\textbf{\ink{1. The price falls when the yield rises}}}

  \vspace{0.2em}
  The result $dV_B/dy<0$, or equivalently $D_{\mathrm{mod}}>0$, follows because
  a higher yield discounts every positive payment more heavily.

  \vspace{0.8em}
  {\large\textbf{\ink{2. A longer par bond is more sensitive to yield}}}

  \vspace{0.2em}
  \[
    D_{\mathrm{mod}}(N)=\frac{1-(1+y/n)^{-N}}{y}
  \]
  increases with $N$. Because each bond starts at $V_B=V_P$, the dollar
  sensitivity $-dV_B/dy=V_PD_{\mathrm{mod}}(N)$ also increases with maturity.
\end{frame}

\begin{frame}{Malkiel's Theorem 3}
  \vspace{0.5em}
  Adding one payment period changes modified duration by:
  \[
    D_{\mathrm{mod}}(N+1)-D_{\mathrm{mod}}(N)
    =\frac{(1+y/n)^{-N-1}}{n}>0.
  \]

  \vspace{0.5em}
  The next increment equals the previous increment multiplied by:
  \[
    (1+y/n)^{-1}<1.
  \]

  Modified duration therefore increases with maturity, but each additional
  period adds less than the previous one. As $N\rightarrow\infty$, modified
  duration approaches $1/y$.

  \slidenote{Limit:}{$1/y$ is the modified duration of a perpetuity.}
\end{frame}

\begin{frame}{Malkiel's Theorem 4}
  \vspace{0.5em}
  Start at the same yield and compare equal changes $\pm\Delta y$.

  \vspace{0.5em}
  \[
    \frac{\Delta V_B}{V_B}\approx
    -D_{\mathrm{mod}}\Delta y+\frac{1}{2}K(\Delta y)^2
  \]

  \vspace{0.5em}
  The duration term changes sign when the direction of the yield change changes.
  The positive convexity term does not. Therefore, a yield decrease raises the
  price more than an equal yield increase lowers it.

  \vspace{0.5em}
  The direction of the curvature matters. A concave price-yield curve would
  produce the opposite asymmetry.
\end{frame}

\begin{frame}{Malkiel's Theorem 5}
  \vspace{0.5em}
  A higher coupon shifts more of the bond's present value to earlier payment
  dates. This lowers Macaulay duration and therefore lowers
  $D_{\mathrm{mod}}$.

  \vspace{0.8em}
  This comparison requires $N\ge2$ and concerns \hilite{percentage
  sensitivity}. Dollar sensitivity can still increase because a higher coupon
  raises the bond's price:
  \[
    -\frac{dV_B}{dy}=V_BD_{\mathrm{mod}}.
  \]

  \slidenote{Edge case:}{at $N=1$, every payment arrives on one date and $D_{\mathrm{mod}}=1/[n(1+y/n)]$, regardless of coupon size.}
\end{frame}

\begin{frame}{Why Theorems 2 and 3 Compare Bonds at Par}
  \vspace{0.3em}
  Away from par, percentage sensitivity and dollar sensitivity can rank bonds
  differently. For a zero-coupon bill,
  $D_{\mathrm{mod}}=T/(1+y/n)$ continues to grow with maturity, but dollar
  sensitivity contains one factor that grows and another that shrinks:
  \[
    -\frac{dV_B}{dy}=V_P\,\underbrace{\frac{T}{1+y/n}}_{\text{\footnotesize growing}}
    \underbrace{\left(1+y/n\right)^{-nT}}_{\text{\footnotesize shrinking}} .
  \]
  With $V_P=100$, $n=1$, and $y=5\%$, dollar sensitivity increases at first,
  reaches a maximum near $T\approx20.5$ years, and then decreases. It is about
  718 USD per unit of yield at the maximum and about 541 at $T=40$.

  \vspace{0.3em}
  The 40-year bill is more sensitive in \hilite{percentage} terms but less
  sensitive in \hilite{dollars}. Starting both bonds at par makes the rankings
  agree.

  \slidenote{Example:}{\dimmed{CHEME-5660-L2b-Example-Treasury-Duration-Convexity-Fall-2026.ipynb}, which tests Theorems 4 and 5.}
\end{frame}

% ---- yield curve ---------------------------------------------
\begin{frame}{Yield Measures and the Yield Curve}
  \vspace{0.3em}
  \begin{itemize}\setlength{\itemsep}{0.45em}
    \item \hilite{Yield to maturity} is the single rate that reproduces a
      security's market price from its promised cash flows.
    \item A \hilite{spot rate} applies to one payment date, so the spot curve
      provides a separate discount factor for each date.
    \item A \hilite{par yield} is the coupon rate that prices a hypothetical bond
      at face value under the spot curve. The published Treasury yield curve is a
      par yield curve.
  \end{itemize}

  Let $D(0,t_j)$ be the time-0 discount factor for a payment at date $t_j$.
  Then:
  \[
    V_B=\sum_{j=1}^{N}X_j\,D(0,t_j).
  \]

  \vspace{0.2em}
  The level, slope, and curvature of the yield curve describe how discounting
  varies with maturity. They do not, by themselves, determine future short
  rates.
\end{frame}

\begin{frame}{Inverted Yield Curves}
  \vspace{0.4em}
  A yield curve is \hilite{inverted} when a long-maturity yield is below a
  short-maturity yield.

  \vspace{0.7em}
  An inverted spread such as \hilite{10-year minus 3-month} has preceded every
  U.S. recession of the past half century.

  \vspace{0.7em}
  This spread does not tell us when a recession will begin. False positives
  occur, lead times range from months to years, and the result depends on the
  maturity pair. The optional par-yield example computes the spread from
  Treasury data.

  \slidenote{Notebook:}{\dimmed{advanced/yield\_curve/CHEME-5660-L3a-ParYieldCurves-Example-Fall-2026.ipynb}}
\end{frame}

\begin{frame}{Optional Advanced Material}
  \vspace{0.2em}
  {\normalsize\textbf{\ink{Malkiel:}} \dimmed{advanced/malkiel/}}

  The Malkiel notebooks study Treasury-bill maturity data and derive a
  closed-form proof of Theorems 2 and 3.

  \vspace{0.5em}
  {\normalsize\textbf{\ink{Yield curve:}} \dimmed{advanced/yield\_curve/}}

  The yield-curve notebooks compute investment yield, solve for yield to maturity
  with the secant method, and construct par yield curves.

  \vspace{0.5em}
  {\normalsize\textbf{\ink{STRIPS:}} \dimmed{advanced/strips/}}

  The STRIPS notebooks infer spot, par, and implied forward rates from Treasury
  STRIPS prices.

  \vspace{0.5em}
  These notebooks are optional and are not prerequisites for L3a.
  \slidenote{Notebooks:}{\href{https://github.com/varnerlab/CHEME-5660-CourseRepository-Fall-2026/tree/main/lectures/week-2/L2b/advanced}{lectures/week-2/L2b/advanced}, also linked from the lecture notebook.}
\end{frame}

% ---- summary -------------------------------------------------
\begin{frame}{Summary}
  We used cash-flow pricing to measure the interest-rate risk of fixed-income
  securities.

  \vspace{0.45em}
  {\large\textbf{\ink{Key takeaways}}}
  \vspace{0.2em}
  \begin{itemize}\setlength{\itemsep}{0.3em}
    \item \textbf{\ink{Price and yield move in opposite directions.}} For fixed
      positive cash flows, the price-yield curve is decreasing and convex.
    \item \textbf{\ink{Duration and convexity estimate small price changes.}}
      The sensitivity depends on the maturity, coupon, and current yield.
    \item \textbf{\ink{Different yield measures answer different questions.}}
      Yield to maturity summarizes one security, spot rates price individual
      payment dates, and par yields summarize the spot curve as coupon rates.
  \end{itemize}
\end{frame}

% ---- closing -------------------------------------------------
\transitionframe{How can \hilite{duration and convexity} help us manage the
interest-rate risk of a bond portfolio?}

\end{document}
